% Class-compatibility preamble retained for a reader excerpt of
%   Dave Witte Morris, "A Cohomological Proof that Real Representations of
%   Semisimple Lie Algebras Have Q-Forms", SIGMA 11 (2015) 034,
%   DOI 10.3842/SIGMA.2015.034, CC BY 4.0.
%
% The source pins the publisher class sigma, which is not part of the delivery
% baseline, so this file supplies the amsart baseline and then copies the author's
% own declaration block mechanically and unchanged:
%   - pinned-original-source.tex lines 11-31: the author's DeclareMathOperator
%     list, the equation numbering and the author's newtheorem declarations;
%   - pinned-original-source.tex lines 35-43, which the author writes immediately
%     after \begin{document} in the publisher layout: the author's own \Res,
%     \ResOne and \Lie macros and \allowdisplaybreaks.
% The publisher header commands of the sigma class (\FirstPageHeading,
% \ArticleName and the like) and the macro \PaperNumber they require are not
% reproduced, because this excerpt carries neither the journal front matter nor
% anything else that uses them.
% No mathematical macro, statement, symbol or assumption is introduced, changed
% or removed by this file.  The author's full DeclareMathOperator list of pinned
% source lines 11-19 is carried, including GL, SL, SO and SU: these four are not
% supplied by the delivery baseline, so leaving them out makes the author's own
% notation (for instance GL_n(Q) in the corollary on lattices) undefined.
\documentclass[10pt,a4paper]{amsart}
\usepackage{amssymb,amsmath,amsthm}
\usepackage{enumerate}
\usepackage{hyperref}

\DeclareMathOperator{\Cliff}{Clif\/f}
\DeclareMathOperator{\Gal}{Gal}
\DeclareMathOperator{\GL}{GL}
\DeclareMathOperator{\Inn}{Inn}
\DeclareMathOperator{\Mat}{Mat}
\DeclareMathOperator{\Out}{Out}
\DeclareMathOperator{\SL}{SL}
\DeclareMathOperator{\SO}{SO}
\DeclareMathOperator{\SU}{\mathbf{SU}}
\numberwithin{equation}{section}

\newtheorem{prop}{Proposition}[section]
\newtheorem{cor}[prop]{Corollary}
\newtheorem{lem}[prop]{Lemma}
{\theoremstyle{definition}
\newtheorem{defn}[prop]{Definition}
\newtheorem{notation}[prop]{Notation}
\newtheorem{rem}[prop]{Remark}
\newtheorem{ques}[prop]{Question}
\newtheorem{eg}[prop]{Example}
\newtheorem{erratum}[prop]{Erratum}}

\newcommand{\Res}[1]{\mathop{\mathbf{R}_{#1}}}
\newcommand{\ResOne}[1]{\mathop{\mathbf{R}^{(1)}_{#1}}}
\newcommand{\Lie}[1]{\mathfrak{\lowercase{#1}}}
\allowdisplaybreaks
